\documentclass{article}

\usepackage{polski}
\usepackage[utf8]{inputenc}

\title{Sprawozdanie z pracowni rotacyjnej}
\author{Jakub Bożek}
\date{\today}

\begin{document}

\maketitle
\tableofcontents
\newpage

\section{Wprowadzenie}
    Na zajęciach z pracowni rotacyjnej u dr. Anity Dąbrowskiej
     chronologia pojęć wyglądała następująco:\\
    \begin{enumerate}
        \item entropia
        \item sposoby liczenia entropii
        \item typy impulsów wzbudzających atom
        \item prawdopodobieństwo wzbudzenia atomu/wystąpienia elektronu dla impulsów
        \item entropia na przykładzie powyższych prawdopodobieństw
        \item \dots
    \end{enumerate}


\section{Entropia}
    Entropia jest to funkcja stanu określająca kierunek przebiegu procesów spontanicznych.
    Jest miarą stopnia nieuporządkowania układu.

    Żeby lepiej zrozumieć czym jest entropia można wziąć na przykład rzut kością 6-ścienną.
    Każdy rzut można określić jako mikrostan, jeśli rzucimy nią 10 razy będziemy mieli jeden makrostan zawierający
    10 mikrostanów.

    Z kombinatoryki wiemy, że aby policzyć ilość możliwych makrostanów należy podnieść liczbę możliwych mikrostanów(6)
    do potęgi określanej przez ilość rzutów(10). $6^{10} = 60,466,176$, tyle wynosi liczba 
\end{document}